% !TEX program = pdflatex
\documentclass[11pt]{article}
\usepackage[margin=1in]{geometry}
\usepackage{amsmath,amsthm,amssymb,mathtools}
\usepackage{hyperref}
\newtheorem{proposition}{Proposition}
\newtheorem{lemma}[proposition]{Lemma}
\newtheorem{definition}{Definition}
\newtheorem{remark}{Remark}
\DeclareMathOperator{\diag}{diag}
\newcommand{\R}{\mathbb R}
\newcommand{\absdet}[1]{\left|\det\!\left(#1\right)\right|}
\title{An Observation-Aware Conditional TT Proposal for Zhao--Cui Filtering}
\author{Mathematical derivation and audit-ready specification}
\date{2026-09-12}

\begin{document}
\maketitle

\begin{abstract}
This note derives an observation-aware coordinate guide for the conditional
tensor-train (TT) proposal used in Zhao--Cui Algorithm~3. The guide changes
the coordinates in which the exact sequential target is fitted; it does not
replace the transition or observation factors, and it does not replace the
conditional TT density by a Gaussian approximation. Propositions establish
the transformed density, its Jacobian, normalization of the conditional
proposal, and exactness of the resulting importance-weight identity.
\end{abstract}

\section{Target and proposal contract}
Let $c=(\theta,x_{t-1})\in\R^{d+m}$ and let $x=x_t\in\R^m$. Given
$y_{1:t}$, define the unnormalized one-step target
\begin{equation}\label{eq:target}
 q_t(x,c)=\widehat\pi_{t-1}(c)\,f_t(x\mid c)\,g_t(y_t\mid x,\theta),
\end{equation}
where $\widehat\pi_{t-1}$ is the previous approximation, $f_t$ is the state
transition density, and $g_t$ is the observation density. Zhao and Cui's
Algorithm~3 draws $x$ conditionally on each retained $c$ and updates the
particle weight with
\begin{equation}\label{eq:weight}
 \omega_t(x,c)=\frac{f_t(x\mid c)g_t(y_t\mid x,\theta)}
 {\widehat p_t(x\mid c,y_{1:t})}.
\end{equation}
The denominator must be the density of the physical state actually drawn.

\section{Observation-dependent affine guide}
\begin{definition}[Affine guide]
For each $c$ and $y_t$, let $m_t(c,y_t)\in\R^m$ and let
$L_t(c,y_t)\in\R^{m\times m}$ be nonsingular. Define
\begin{equation}\label{eq:guide}
 G_t(u;c,y_t)=m_t(c,y_t)+L_t(c,y_t)u,\qquad u\in\R^m.
\end{equation}
The guide may be obtained from weighted particle moments, a standard UKF
moment calculation, or a likelihood-weighted sigma-point projection. Those
moment sources are distinct methods; the derivation below only requires the
nonsingularity of $L_t$.
\end{definition}

\begin{proposition}[Pullback and Jacobian]
\label{prop:pullback}
For fixed $(c,y_t)$, the map $u\mapsto x=G_t(u;c,y_t)$ is a bijection and
\begin{equation}\label{eq:pullback}
 q_t^G(u,c)=q_t(G_t(u;c,y_t),c)\,\absdet{L_t(c,y_t)}.
\end{equation}
If the joint map is $\widetilde G_t(u,c)=(G_t(u;c,y_t),c)$, its Jacobian
determinant is also $\det L_t(c,y_t)$.
\end{proposition}
\begin{proof}
Nonsingularity gives the inverse
$G_t^{-1}(x;c,y_t)=L_t(c,y_t)^{-1}(x-m_t(c,y_t))$. The ordinary change-of-
variables formula gives~\eqref{eq:pullback}. The Jacobian of $\widetilde G_t$
has block form
\[
 \begin{pmatrix}L_t&\partial_cG_t\\0&I\end{pmatrix},
\]
which is block triangular, so its determinant is $\det L_t$.
\end{proof}

\begin{proposition}[Conditional TT proposal]
\label{prop:proposal}
Suppose a nonnegative squared-TT approximation produces a normalized
conditional density $\widehat p_t^G(u\mid c,y_{1:t})$. Define the physical
proposal by pushing it through~\eqref{eq:guide}:
\begin{equation}\label{eq:physical-proposal}
 \widehat p_t(x\mid c,y_{1:t})
 =\widehat p_t^G\!\left(G_t^{-1}(x;c,y_t)\mid c,y_{1:t}\right)
   \absdet{L_t(c,y_t)}^{-1}.
\end{equation}
Then $\widehat p_t(\cdot\mid c,y_{1:t})$ is normalized and sampling
$u\sim\widehat p_t^G$ followed by $x=G_t(u;c,y_t)$ has this density.
\end{proposition}
\begin{proof}
Substitute $u=G_t^{-1}(x;c,y_t)$ in the change-of-variables formula. The
factor $\absdet{L_t}^{-1}$ is the inverse Jacobian. Integrating~\eqref{eq:physical-proposal}
over $x$ and changing variables back to $u$ gives integral one.
\end{proof}

\section{How the TT fit uses the observation}
The TT regression target in guide coordinates is the exact pullback~\eqref{eq:pullback}.
Equivalently, following Zhao--Cui's bridge construction, choose a positive
bridge $\rho_t$ and a reference density $\eta$ for $u,c$, then fit
\begin{equation}\label{eq:ratio}
 q_{t,\rho}^G(u,c)=
 \frac{q_t(G_t(u;c,y_t),c)}{\rho_t(G_t(u;c,y_t),c)}\,\eta(u,c).
\end{equation}
The fitted squared TT and its conditional KR rearrangement operate on this
transformed object. The observation enters through $g_t$ in~\eqref{eq:target}
and, when the guide is adaptive, through $m_t,L_t$ and the inverse map.

\begin{proposition}[Exact importance correction]
\label{prop:weights}
Assume the previous approximation is the factor appearing in~\eqref{eq:target}
and that the proposal is~\eqref{eq:physical-proposal}. Then the Algorithm~3
increment~\eqref{eq:weight} is the Radon--Nikodym correction from the proposal
joint density to the target joint density. In particular, omitting the
Jacobian in~\eqref{eq:physical-proposal}, or dividing by a bare guide Gaussian,
does not give the stated correction unless that omitted factor is identically
one.
\end{proposition}
\begin{proof}
The sampling density for $(c,x)$ is
$\widehat\pi_{t-1}(c)\widehat p_t(x\mid c,y_{1:t})$. Dividing~\eqref{eq:target}
by this density cancels $\widehat\pi_{t-1}$ and yields~\eqref{eq:weight}.
Substitution of~\eqref{eq:physical-proposal} shows that the inverse Jacobian
must be present in the denominator. Removing it changes the Radon--Nikodym
derivative by $\absdet{L_t}$.
\end{proof}

\begin{proposition}[Observation-response requirement]
\label{prop:response}
If $m_t$ and $L_t$ are independent of $y_t$, then the guide itself cannot
provide observation adaptation. If, additionally, the TT rows and fitted
conditional are unchanged when $y_t$ changes, the resulting proposal is not
observation-responsive even though the exact likelihood remains in the final
weight.
\end{proposition}
\begin{proof}
The map~\eqref{eq:guide} depends on the observation only through $m_t,L_t$.
If those and the fitted conditional are unchanged, the distribution of the
generated $x$ conditional on $c$ is unchanged. The likelihood can still alter
the importance weight, but it has not altered the proposal.
\end{proof}

\section{Method status and audit boundary}
The source-faithful baseline is Zhao--Cui's Gaussian bridge: weighted particle
moments define a full-joint affine map, followed by fitting the transformed
target-to-bridge ratio. Their tempered bridge is the source-supported nonlinear
extension. A standard UKF is a comparator for constructing moments; it uses
transformed observation sigma points and state--observation cross-covariance.
Likelihood-softmax weighting of state sigma points is a separate extension.

The present derivation proves the mathematical contract conditional on a
nonsingular guide and a normalized TT conditional. It does not prove that UKF
moments, the likelihood-weighted projection, or any finite TT rank gives a
better filter. Those claims require executable response, round-trip, Jacobian,
and replicated filtering tests.

\begin{thebibliography}{9}
\bibitem{zhaocui2024} Y. Zhao and T. Cui, ``Tensor-train methods for sequential
state and parameter learning in state-space models,'' \emph{JMLR} 25 (2024),
1--73. \href{https://www.jmlr.org/papers/v25/23-0743.html}{Official article}.
\bibitem{julieruhlmann1997} S. J. Julier and J. K. Uhlmann, ``A new extension
of the Kalman filter to nonlinear systems,'' Proc. SPIE 3068 (1997), 182--193,
doi:10.1117/12.280797.
\bibitem{cuidolgov2022} T. Cui and M. Dolgov, ``Deep composition of
tensor-trains using squared inverse Rosenblatt transports,'' \emph{Foundations
of Computational Mathematics} 22 (2022), 1261--1296,
doi:10.1007/s10208-021-09537-5.
\end{thebibliography}
\end{document}
